%   Packages and macros of the recovered 2024 edition.

\usepackage{amsmath}
\usepackage{array}
\usepackage{longtable}
\usepackage{amssymb}
\usepackage{amsthm}
\usepackage{graphicx}
\usepackage{subcaption}
\usepackage{algorithm}
\usepackage{algpseudocode}
\usepackage{tikz}
\usetikzlibrary{calc}
\usepackage[edges]{forest}
\usepackage{xcolor}
\usepackage{hyperref}

\frenchspacing
\renewcommand*{\chapterautorefname}{Chapter}
\renewcommand*{\sectionautorefname}{Section}
\renewcommand*{\subsectionautorefname}{Section}

% Theorem-like environments: one counter, numbered within sections.
\theoremstyle{definition}
\newtheorem{definition}{Definition}[section]
\theoremstyle{plain}
\newtheorem{proposition}[definition]{Proposition}
\newtheorem{lemma}[definition]{Lemma}
\newtheorem{corollary}[definition]{Corollary}
\newtheorem{theorem}[definition]{Theorem}
\newtheorem{claim}[definition]{Claim}
\newtheorem{observation}[definition]{Observation}
\newtheorem{problem}[definition]{Problem}
\theoremstyle{remark}
\newtheorem{example}[definition]{Example}

% Names set in typewriter type in the deposited PDF.
\newcommand{\FKA}{\texttt{FK-A}}
\newcommand{\FKB}{\texttt{FK-B}}
\newcommand{\TransHyp}{\texttt{Trans-Hyp}}
\newcommand{\Tr}{\mathit{Tr}}
\newcommand{\Aut}{\mathit{Aut}}
\newcommand{\GK}{\mathcal{GK}}
\newcommand{\TK}{\mathcal{TK}}

% Automorphism-tree nodes (Figures 5.1-5.4 and the appendix).
\newcommand{\anode}[5]{\textbf{Node: #1}\\Split vert.: #2\\Aut Grp: #3\\Hyp.: #4\\Graph: #5}
\forestset{
  autotree/.style={
    for tree={draw, align=center, font=\tiny, inner sep=1.5pt,
      forked edge, l sep=2.5mm, s sep=1mm, fork sep=1.2mm, anchor=north,
      child anchor=north, parent anchor=south},
  },
}

% Hypergraph drawings (generated by ../tools/hyperfig.py).
\tikzset{
  hgfill/.style={fill=#1, fill opacity=0.35, draw=none},
  hgline/.style={draw=black, line width=0.9pt, fill=none},
  hgvertex/.style={inner sep=0pt, font=\normalsize},
}
\hypersetup{hidelinks}

% Name of the Lean declaration proving a statement (revised edition).
\newcommand{\leanref}[1]{{\footnotesize\textsf{[Lean:~\texttt{\detokenize{#1}}]}}}
